
\section{De una lámina a un laminado: qué se suma}
Una lámina ortótropa tiene rigideces distintas según dirección y acoplamiento por Poisson. Al girar la lámina se transforma su matriz de rigidez; no se promedian simplemente módulos longitudinal y transversal. En teoría clásica de laminados, bajo sus hipótesis de placa delgada y unión entre capas, resultantes y deformaciones se relacionan mediante
\begin{equation}\begin{bmatrix}\mathbf N\\\mathbf M\end{bmatrix}=\begin{bmatrix}\mathbf A&\mathbf B\\\mathbf B&\mathbf D\end{bmatrix}\begin{bmatrix}\boldsymbol\varepsilon^0\\\boldsymbol\kappa\end{bmatrix}.\end{equation}
$\mathbf A$ reúne rigidez de membrana, $\mathbf D$ la de flexión y $\mathbf B$ el acoplamiento entre ambas. Las unidades son, respectivamente, fuerza por longitud, fuerza por longitud multiplicada por longitud al cuadrado, y fuerza, si se usan las resultantes habituales. La convención de signos de curvatura y del eje de espesor se declara antes de comparar matrices.

Un apilado simétrico respecto al plano medio anula $\mathbf B$ en el modelo ideal porque las contribuciones impares en espesor se cancelan. Eso no elimina todos los acoplamientos posibles dentro de $\mathbf A$ o $\mathbf D$ ni demuestra isotropía. Cambiar el orden de capas puede conservar una rigidez de membrana y modificar flexión. Las propiedades y espesores reales siguen siendo necesarios.
\section{Cargas de matriz, interfaces y daño}
Las fibras pueden dominar una rigidez longitudinal, pero la matriz transmite carga entre fibras, estabiliza y participa en respuesta transversal y cortante. La interfaz determina transferencia y ciertos modos de daño. Bajo compresión, flexión o carga fuera de plano, afirmar que la matriz «no soporta carga» puede ocultar el mecanismo que gobierna.

Fisuración de matriz, rotura de fibra y delaminación pueden aparecer en secuencia e interactuar. Un impacto con huella pequeña puede dejar daño interno relevante; observar la superficie no caracteriza todo el volumen. La resistencia interlaminar y la tenacidad de delaminación describen aspectos distintos y dependen del modo de apertura o cortante. La teoría clásica elástica de placa no resuelve por sí sola el campo tridimensional de un borde libre o un impacto.
\Needspace{.48\textheight}
\section{Sándwich, GLARE y uniones adhesivas}
Separar caras aumenta rigidez a flexión con poca cantidad de material resistente de cara, pero el núcleo debe transmitir cortante y mantener la separación. Arrugamiento, aplastamiento, despegue e infiltración pueden gobernar antes de alcanzar la resistencia nominal de las caras. La ganancia ideal de inercia del problema P25 debe leerse junto a estas limitaciones.

GLARE combina capas metálicas y de compuesto de vidrio en una arquitectura específica; no es cualquier panel con vidrio y aluminio. Las rutas de daño, propagación y reparación dependen de esa arquitectura. En una unión adhesiva, $F/(bL)$ solo es cortante media: extremos del solape y excentricidad introducen concentraciones y pelado. Preparación superficial, curado, espesor, ambiente y compatibilidad de sustratos son parte del sistema resistente.
\repaso{¿Un laminado simétrico y equilibrado es necesariamente isótropo?}{No. Simetría y equilibrio eliminan ciertos acoplamientos bajo condiciones del modelo, pero las rigideces pueden seguir variando con dirección. La isotropía o cuasiisotropía requiere condiciones adicionales sobre el apilado y el comportamiento considerado.}
